% Part: computability
% Chapter: computability-theory
% Section: applications-fixed-point

\documentclass[../../../include/open-logic-section]{subfiles}

\begin{document}

\olfileid{cmp}{thy}{apf}
\olsection{స్థిరబిందు సిద్ధాంతపు అనువర్తనం}

పాక్షిక గణనీయ ప్రమేయాలను వాటి సూచికల పరంగా నిర్వచించడానికి
స్థిరబిందు సిద్ధాంతం ముఖ్యంగా వీలు కల్పిస్తుంది. ఉదాహరణకు, ప్రతి~$y$కు
\[
\cfind{e}(y)=e+y.
\]
అయ్యే సూచిక~$e$ను కనుగొనవచ్చు. మరొక ఉదాహరణగా, స్థిరబిందు సిద్ధాంతపు
నిరూపణను వాడి తననే తాను ముద్రించే Java లేదా C++ ప్రోగ్రామును
రూపొందించవచ్చు.

ప్రతి~$e$కు $W_e$ అనేది $\cfind{e}$ యొక్క నిర్వచన క్షేత్రం అని
ఉంచితే, $W_0$, $W_1$, $W_2$,~\dots{} అనుక్రమం గణనీయంగా
లెక్కించదగిన సమితులను లెక్కిస్తుందని గుర్తుంచుకోండి. వీటిలో కొన్ని
సమితులు గణనీయాలు. ఒక విలువ~$x$ను ఇన్‌పుట్‌గా తీసుకుని, $W_x$
గణనీయం కావడం జరిగితే దాని లక్షణ ప్రమేయానికి ఒక సూచికను తిరిగి ఇచ్చే
అల్గారిథం ఉందా అని అడగవచ్చు. సమాధానం ``లేదు''; అలాంటి అల్గారిథం లేదు:

\begin{thm}
కింది ధర్మం గల పాక్షిక గణనీయ ప్రమేయం~$f$ ఏదీ లేదు: $W_e$
గణనీయమైనప్పుడల్లా $f(e)$ నిర్వచితమై, $\cfind{f(e)}$ అనేది దాని
లక్షణ ప్రమేయం అవుతుంది.
\end{thm}

\begin{proof}
$f$ ఏ పాక్షిక గణనీయ ప్రమేయమైనా అనుకుందాం. $W_e$ గణనీయమై ఉన్నా,
$\cfind{f(e)}$ దాని లక్షణ ప్రమేయం కాని విధంగా ఒక~$e$ను నిర్మిస్తాం.
స్థిరబిందు సిద్ధాంతాన్ని వాడి
\[
\cfind{e}(y) \simeq
\begin{cases}
0 & \text{$y=0$ మరియు $\cfind{f(e)}(0) \fdefined = 0$ అయితే} \\
\text{అనిర్వచితం} & \text{ఇతర సందర్భాల్లో}
\end{cases}
\]
అయ్యే సూచిక~$e$ను కనుగొనవచ్చు. అంటే, కింది విధంగా నిర్వచించిన
ప్రమేయానికి స్థిరబిందు సిద్ధాంతాన్ని వర్తింపజేసి $e$ను పొందుతాం:
\[
g(x,y) \simeq
\begin{cases}
0 & \text{$y=0$ మరియు $\cfind{f(x)}(0) \fdefined = 0$ అయితే} \\
\text{అనిర్వచితం} & \text{ఇతర సందర్భాల్లో.}
\end{cases}
\]
$g$ పాక్షిక గణనీయమని అనౌపచారికంగా ఇలా చూడవచ్చు: ఇన్‌పుట్‌లు
$x$,~$y$పై, అల్గారిథం ముందుగా $y$ అనేది~$0$కు సమానమో లేదో
తనిఖీ చేస్తుంది. సమానమైతే అది $f(x)$ను గణించి, తరువాత సార్వత్రిక
యంత్రాన్ని వాడి $\cfind{f(x)}(0)$ను గణిస్తుంది. ఈ చివరి గణన ఆగి
విలువ~$0$ను ఇస్తే అల్గారిథం~$0$ను తిరిగి ఇస్తుంది; లేకపోతే అది ఆగదు.
\sourcecorrection{OLTECOMTHY-018}{సిద్ధాంతం పాక్షిక గణనీయ $f$ ఏదీ ఉండదని చెబుతుంది; కానీ ఆంగ్ల మూల నిరూపణ ``Let $f$ be any computable function'' అని సర్వనిర్వచిత $f$కు మాత్రమే వాదనను మొదలుపెట్టింది. తరువాతి $g$ నిర్మాణం పాక్షిక $f$కూ పనిచేస్తుంది; సిద్ధాంతానికి కావలసినట్లుగా ``పాక్షిక గణనీయ ప్రమేయం''గా సరిచేశాం.}

ఇప్పుడు $\cfind{f(e)}(0)$ నిర్వచితమై $0$కు సమానమైతే,
$\cfind{e}(y)$ అనేది $y$ అనేది~$0$కు సమానమైనప్పుడు మాత్రమే
నిర్వచితమవుతుంది;
కాబట్టి $W_e=\{0\}$. $\cfind{f(e)}(0)$ అనిర్వచితమైతే, లేదా
నిర్వచితమైనా $0$కు సమానం కాకపోతే, $W_e=\emptyset$. ఏ సందర్భంలోనైనా
$\cfind{f(e)}$ అనేది $W_e$ యొక్క లక్షణ ప్రమేయం కాదు; ఎందుకంటే అది
ఇన్‌పుట్~$0$పై తప్పు సమాధానం ఇస్తుంది.
\end{proof}

\end{document}
